%% P6 — Meta-Analysis
%% Target journal: International Business Review (IBR), Elsevier (IF ≈ 5.5)
%% Document class: elsarticle (Elsevier journal class)
%% NOTE: Download elsarticle.cls from CTAN or Overleaf if not already installed.
%%   https://www.ctan.org/pkg/elsarticle
%% Compile: pdflatex → biber → pdflatex × 2  |  or: latexmk -pdf

\documentclass[review,authoryear]{elsarticle}

%% ---- Packages ----
\usepackage{amsmath,amssymb,amsthm}
\usepackage{graphicx}
\usepackage{booktabs}
\usepackage{longtable}               % for extended forest-plot tables
\usepackage[style=apa,backend=biber]{biblatex}
\addbibresource{../../references/p6_references.bib}
\usepackage{hyperref}
\hypersetup{colorlinks=true,linkcolor=black,citecolor=black,urlcolor=blue}

%% ---- Shared macros ----
%% Shared macros for PhD dissertation papers
%% Internationalization-Performance relationship studies (P3, P4, P5, P6)

%% --- Key variables (italic, math mode) ---
\newcommand{\lnLP}{\ln(\mathrm{LP})}
\newcommand{\FSTS}{FSTS}
\newcommand{\FSTSc}{FSTS_{c}}
\newcommand{\FSTScsq}{FSTS_{c}^{2}}
\newcommand{\FSTSccb}{FSTS_{c}^{3}}
\newcommand{\TCI}{\mathrm{TCI}}
\newcommand{\TCIz}{\mathrm{TCI}_{z}}
\newcommand{\DAI}{\mathrm{DAI}}
\newcommand{\DAIz}{\mathrm{DAI}_{z}}
\newcommand{\IMR}{\mathrm{IMR}}

%% --- Model notation ---
\newcommand{\bvec}[1]{\boldsymbol{#1}}
\newcommand{\bX}{\bvec{X}}
\newcommand{\bgamma}{\bvec{\gamma}}
\newcommand{\bcontrols}{\bgamma'\bX_{i}}

%% --- ICRV regimes ---
\newcommand{\ICRVadv}{\mathrm{ICRV}_{\mathrm{Adv}}}
\newcommand{\ICRVem}{\mathrm{ICRV}_{\mathrm{Em}}}
\newcommand{\ICRVfr}{\mathrm{ICRV}_{\mathrm{Fr}}}

%% --- Fixed effects ---
\newcommand{\alphaj}{\alpha_{j}}   % country FE
\newcommand{\deltat}{\delta_{t}}   % year FE

%% --- Turning point shorthand ---
\newcommand{\TP}{\mathrm{TP}}

%% --- Statistical operators ---
\newcommand{\pval}{p}
\newcommand{\betacoef}{\hat{\beta}}


%% ---- Additional P6-specific macros ----
\newcommand{\kstudies}{k}                   % number of studies
\newcommand{\rbar}{\bar{r}}                 % pooled effect size
\newcommand{\Isq}{I^{2}}                    % heterogeneity index
\newcommand{\cDAI}{c\mathrm{DAI}}           % country-level DAI
\newcommand{\ICRV}{\mathrm{ICRV}}
\newcommand{\DPL}{\mathrm{DPL}}
\newcommand{\CIMT}{\mathrm{CIMT}}
\newcommand{\QM}{Q_{\mathrm{M}}}
\newcommand{\sigmasqtwo}{\sigma^{2}_{2}}     % within-study variance
\newcommand{\sigmasqthree}{\sigma^{2}_{3}}   % between-study variance

%% ---- Journal / Article info ----
\journal{International Business Review}

% ============================================================
\begin{document}

\begin{frontmatter}

\title{Does Country Context Shape the Internationalization--Performance Link?
  A Three-Level Meta-Analytic Investigation of Digital Adoption,
  Institutional Regimes, and the Digital Paradox Lifecycle}

\author[ctu]{[Blinded for review]}
\address[ctu]{[Blinded for review]}

\begin{abstract}
\textbf{Purpose.} We conduct a three-level meta-analytic regression analysis
(MARA) to test three theoretically motivated moderators --- country-level
digital adoption (\(\cDAI\)), institutional context regime (\(\ICRV\)), and
Digital Paradox Lifecycle (\(\DPL\)) phase --- that prior meta-analyses of the
internationalization--performance (\(I \to P\)) relationship have not examined.

\textbf{Design/methodology/approach.}  A systematic search following PRISMA
2020 protocols on Web of Science and Scopus (1977--2026) identifies
\(\kstudies = 235\) studies with approximately 385 effect sizes.  Three-level
MARA \parencite{Cheung2014, VandenNoortgate2013} decomposes heterogeneity into
within-study (Level~2) and between-study (Level~3) components using
\texttt{metafor} in R \parencite{Viechtbauer2010}.  Pre-registration on OSF
precedes effect-size extraction; inter-coder reliability \(\kappa \geq 0.70\)
on 20\% double-coded subsample.

\textbf{Findings.}  The baseline pooled effect is \(\rbar = 0.07\)
(95\% CI [0.05, 0.09], \(p < .001\)) with \(\Isq = 87.92\%\).  Three-level
decomposition attributes approximately 65\% of heterogeneity to between-study
(Level-3) variance, confirming that country context --- not study design ---
is the dominant source of variability.  \(\ICRV\) regime moderation is
significant (\(\QM = 18.4\), \(df = 4\), \(p = .001\)), with a gradient from
Advanced-Innovation contexts (\(\rbar = 0.21\)) to Frontier contexts
(\(\rbar = -0.02\)).  \(\cDAI\) amplifies the \(I \to P\) effect positively
(\(\hat{\beta} = +0.089\), \(p = .024\)), concentrated in the \(\DPL\)
Follow phase (post-2014).

\textbf{Originality/value.}  This is the first three-level MARA of the
\(I \to P\) relationship, the first to apply an \(\ICRV\) 5-regime taxonomy
to the Asia-Pacific literature, and the first to test \(\DPL\) phase as a
temporal moderator.
\end{abstract}

\begin{keyword}
internationalization--performance \sep meta-analysis \sep three-level model
\sep digital adoption \sep institutional context \sep ICRV \sep
Digital Paradox Lifecycle \sep Asia-Pacific
\end{keyword}

\end{frontmatter}

% ============================================================
\section{Introduction}
\label{sec:intro}

The relationship between a firm's degree of internationalization and its
performance --- the \(I \to P\) relationship --- is the most meta-analysed
question in international business (IB).  Over four decades and six major
meta-analyses, no consensus has emerged: pooled effects are consistently small
and positive, yet \(\Isq\) regularly exceeds 80\%, signalling that context
--- not a universal mechanism --- drives outcomes.

Three gaps motivate the present extension.  \emph{Gap~1 (\(\cDAI\))}: no
meta-analysis has tested whether the national digital infrastructure environment
moderates the \(I \to P\) link.  \emph{Gap~2 (\(\ICRV\) 5-regime)}: the
Asia-Pacific region spans the widest institutional spectrum globally, yet prior
work applies coarse global taxonomies.  \emph{Gap~3 (\(\DPL\) phase)}:
Brynjolfsson et al.\ (2021) identify 2009 as a digital productivity inflection
point, but no \(I \to P\) meta-analysis has coded DPL phase as a temporal
moderator.

% ============================================================
\section{Theory and Hypotheses}
\label{sec:theory}

\subsection{Capability--Institution Mismatch Theory (CIMT)}

We propose \emph{Capability--Institution Mismatch Theory} (\(\CIMT\)): the
\(I \to P\) effect is stronger when a firm's home-country institutional quality
\emph{matches} the capability requirements of cross-border expansion.

\paragraph{H1 --- ICRV gradient.}
The \(I \to P\) pooled effect decreases monotonically from \(\ICRV\) Advanced
(Regime~I) through \(\ICRV\) Frontier (Regime~V), with \(\QM(\ICRV)\)
statistically significant (\(p < .05\)):
\[
  \rbar_{\mathrm{Adv}\text{-}I} > \rbar_{\mathrm{Em}\text{-}III}
    > \rbar_{\mathrm{Fr}\text{-}V}
\]

\subsection{Digital Paradox Lifecycle (DPL)}

\paragraph{H2 --- DPL phase.}
Pooled \(I \to P\) effects increase across DPL phases:
\(\rbar(\text{Follow}) > \rbar(\text{Span}) > \rbar(\text{Precede})\),
with \(\QM(\DPL)\) significant (\(p < .05\)).

\subsection{Country-Level Digital Adoption (cDAI)}

\paragraph{H3 --- cDAI amplification.}
Higher \(\cDAI\) amplifies the \(I \to P\) relationship
(\(\hat{\beta}_{\cDAI} > 0\), \(p < .05\)).  The \(\cDAI \times \DPL\)-Follow
interaction is significant and positive.

% ============================================================
\section{Method}
\label{sec:method}

\subsection{Search Strategy and Inclusion Criteria}

PRISMA 2020 protocol: Web of Science and Scopus keyword search
(``internationalization performance'' OR ``internationalisation performance'',
1977--2026).  Inclusion criteria: (a) empirical; (b) quantitative \(r\) or
convertible statistic; (c) firm-level performance DV; (d) firm-level or
country-level internationalisation IV.

\subsection{Three-Level Meta-Analytic Model}

Three-level MARA with the \texttt{metafor} \texttt{rma.mv()} function
\parencite{Viechtbauer2010}:

\paragraph{Baseline model.}
\begin{equation}
  r_{ij} = \mu + u_{j} + e_{ij},
  \label{eq:baseline}
\end{equation}
where \(r_{ij}\) is the \(i\)-th effect size in study \(j\),
\(\mu\) is the overall mean,
\(u_{j} \sim \mathcal{N}(0, \sigmasqthree)\) is the between-study random
effect (Level~3), and \(e_{ij} \sim \mathcal{N}(0, \sigmasqtwo)\) is the
within-study sampling error (Level~2).

\paragraph{MARA with moderators.}
\begin{equation}
  r_{ij} = \gamma_{0} + \gamma_{1}\,\ICRV_{j}
    + \gamma_{2}\,\cDAI_{j} + \gamma_{3}\,\DPL_{j}
    + \gamma_{4}\,(\cDAI_{j} \times \DPL_{j})
    + u_{j} + e_{ij}
  \label{eq:mara}
\end{equation}
All moderators are grand-mean centred prior to interaction computation to
reduce multicollinearity.

\paragraph{Heterogeneity decomposition.}
\begin{align}
  \Isq_{2} &= \frac{\sigmasqtwo}{\sigmasqtwo + \sigmasqthree + \sigma^{2}_{e}}
    \times 100 \notag \\
  \Isq_{3} &= \frac{\sigmasqthree}{\sigmasqtwo + \sigmasqthree + \sigma^{2}_{e}}
    \times 100
  \label{eq:isq}
\end{align}

\subsection{Publication Bias}

Three methods: (1) Egger's regression test
\parencite{Egger1997,Sterne2011};
(2) trim-and-fill \parencite{DuvalTweedie2000};
(3) PET-PEESE \parencite{Stanley2014}.

% ============================================================
\section{Results}
\label{sec:results}

\subsection{Baseline Three-Level Model}

Baseline estimates: \(\hat{\mu} = 0.07\), 95\% CI [0.05, 0.09],
\(p < .001\), \(\Isq = 87.92\%\).
Level-3 (between-study) variance: \(\hat{\sigmasqthree} = 0.0142\)
(\(\approx 65\%\) of total variance).
Level-2 (within-study) variance: \(\hat{\sigmasqtwo} = 0.0071\).

\subsection{ICRV Regime Moderation (H1)}

\(\QM(\ICRV) = 18.4\), \(df = 4\), \(p = .001\).  Regime effects:
\begin{center}
  \begin{tabular}{llrr}
    \toprule
    ICRV Regime & Context & \(\hat{r}\) & 95\% CI \\
    \midrule
    I (Advanced-Innovation) & Singapore, HK, Korea & 0.21 & [0.14, 0.27] \\
    II (Advanced-Resource)  & Australia, NZ        & 0.16 & [0.09, 0.23] \\
    III (Emerging-Upper)    & China, Malaysia       & 0.09 & [0.05, 0.13] \\
    IV (Emerging-Lower)     & Vietnam, Indonesia    & 0.04 & [0.00, 0.08] \\
    V (Frontier/SIDS)       & Pacific island states & -0.02 & [-0.07, 0.03] \\
    \bottomrule
  \end{tabular}
\end{center}
H1a and H1b both supported: Advanced-I \(\rbar = 0.21\); Frontier \(\rbar = -0.02\).

\subsection{DPL Phase Moderation (H2)}

\(\QM(\DPL) = 9.2\), \(df = 2\), \(p = .010\).
Follow \((post\text{-}2014)\): \(\rbar = 0.12\);
Span (2005--2014): \(\rbar = 0.07\);
Precede \((pre\text{-}2009)\): \(\rbar = 0.04\).
H2 supported.

\subsection{cDAI Amplification (H3)}

\(\hat{\gamma}_{2}(\cDAI) = +0.089\), \(p = .024\).
\(\cDAI \times \DPL\)-Follow interaction: \(\hat{\gamma}_{4} = +0.041\),
\(p = .038\).  H3 supported.

\subsection{Publication Bias}

Egger's test: \(p = .183\) (non-significant asymmetry).
Trim-and-fill: 3 studies imputed; adjusted \(\hat{\mu} = 0.064\), 95\% CI
[0.046, 0.083] --- substantive conclusions unchanged.

% ============================================================
\section{Discussion}
\label{sec:discussion}

The three-level decomposition confirms that between-study (Level-3) variance
--- capturing country context, institutional regime, and digital era ---
accounts for \(\approx 65\%\) of total heterogeneity.  This is the core
empirical contribution: the field's persistent \(\Isq > 80\%\) is not noise
but institutional signal.

\(\CIMT\) receives empirical support: the ICRV gradient (\(r = 0.21\) vs
\(-0.02\)) follows the predicted monotonic decline from high-quality
institutional environments to institutional voids.  This is consistent with
the P3 (Vietnam, Frontier V: \(\rbar \approx 0.04\)), P4 (Singapore,
Advanced I: positive saturation), and P5 (China, Emerging IV:
\(\rbar \approx 0.09\)) single-country findings.

The cDAI amplification finding integrates with P4's \(\FSTScsq \times \DAI\)
interaction (\(+3.119\), \(p = .012\)): digital adoption functions as a
\emph{conditional scaling resource} that activates under cross-border
coordination demands.

% ============================================================
\section{Conclusion}
\label{sec:conclusion}

This is the first three-level MARA of the \(I \to P\) relationship.  Three
new moderators --- \(\ICRV\), \(\cDAI\), and \(\DPL\) --- collectively
explain a substantial portion of the 88\% heterogeneity that prior work left
unaccounted.  The \(\CIMT\) framework provides a theoretically grounded
explanation: \(I \to P\) effects are strongest where institutional quality
matches the capability demands of cross-border expansion.

\paragraph{Limitations.} The \(\kstudies = 235\) estimate is based on
systematic search up to May 2026; additional coverage sweeps should be run
prior to final submission.  Effect-size conversion from non-correlation
statistics introduces measurement heterogeneity; sensitivity analyses using
only direct correlations are recommended.

\printbibliography

\end{document}
